\subsection{Vertex Minors}

\begin{definition}
A graph $H$ is a vertex minor of a graph $G$ if there is a sequence of vertex deletions and local complementations from $G$ to $H$.
\end{definition}

\note{The above definition is a major motivation for our focus on vertex deletion in the previous section. Thus, it may make sense to swap the order of sections around}

\begin{theorem}
Deciding whether a graph $H$ is a vertex-minor of $G$ is $\NP$-complete \cite{vertex-minor-np}.
\end{theorem}

\begin{definition}
\note{The definition of the \term{rank width} of a graph $G = (V,E)$ is quite involved and given in \cite{rank-width-def}.}
\end{definition}

\begin{fact}
The rank width of a graph is preserved under local complementation. Vertex deletions may reduce rank-width, so the rank width of $G$ is at least the rank width of any vertex-minor of $G$.
\end{fact}

\begin{theorem}
Every $\CGrid$ is a vertex-minor of a graph in every family of graphs with unbounded rank-width [cite:need]. Formally, there is a function $f: \N \to \N$ such that every graph of rank-width $\geq f(n)$ has a vertex-minor isomorphic to $\CGrid_{n \times n}$ \cite{grid-theorem}.
\end{theorem}

\begin{theorem}
Grid graphs contain all graphs as vertex minors. Formally, there is a function $f: G \to \N$ such a graph $G$ is the vertex minor of $\Grid_{f(G) \times f(G)}$.
\begin{proof}
\note{I would like to formalize the proof in Kunal's notes from James George Davies.}
\end{proof}
\end{theorem}

\begin{theorem}
If $\CGrid \geq_\LOCC \Grid$ then unbounded rank-width is not sufficient for universality.
\end{theorem}